\documentclass[11pt]{article}
\usepackage[margin=1in]{geometry}
\usepackage{amsmath,amssymb,booktabs,tabularx}
\usepackage[T1]{fontenc}
\usepackage{lmodern}
\usepackage[hidelinks]{hyperref}
\usepackage{listings}
\lstset{basicstyle=\small\ttfamily,breaklines=true,columns=fullflexible,
  keepspaces=true,frame=single}
\setlength{\parindent}{0pt}
\setlength{\parskip}{6pt}
\setlength{\emergencystretch}{3em}
\newcommand{\vect}[1]{\boldsymbol{#1}}
\newcommand{\R}{\mathbb{R}}

\title{Neo-Hookean Liver Deformation:\\Model, Numerical Solution, and Implementation Outline}
\author{}
\date{18 September 2026}

\begin{document}
\maketitle

\textbf{Scope.} This document describes the deformation problem implemented in
\path{simulate_deformation/}, including its current Elmer interface. It then
outlines a custom solver for the same boundary-value problem. The custom solver
and incremental-loading procedure are proposed designs, not implemented features
of the current generation script. No voxelization or registration is required
to generate a deformation.

\section*{Example outline}
\begin{enumerate}
  \item Input mesh, units, and unknown displacement field.
  \item Fixed boundary and traction-patch sampling.
  \item Neo-Hookean material and equilibrium equations.
  \item Tetrahedral discretization and force assembly.
  \item Nonlinear updates, linear solves, and line search.
  \item Convergence criteria, failure classification, and output validation.
  \item Current file workflow and reproducible settings.
  \item Implementation outline for a custom solver.
\end{enumerate}

\section{Input and unknowns}
We start with a volumetric liver mesh consisting of $N$ vertices
$\vect{X}_a\in\R^3$ and four-node tetrahedra. Input coordinates to
\path{main.py} are in metres. Young's modulus is in pascals, traction is in
pascals, and forces are in newtons. The MSD batch wrapper converts source
coordinates from millimetres to metres before invoking this script.

The unknown is the displacement $\vect{u}_a$ at every vertex. The deformed
position is $\vect{x}_a=\vect{X}_a+\vect{u}_a$. We solve a static equilibrium
problem without inertia or body forces. The reference configuration is
$\Omega_0$; its boundary is divided into a fixed region $\Gamma_D$, loaded
patches $\Gamma_{p,0}$, and a traction-free remainder.

\section{Boundary conditions and sampling}
Exterior triangles are extracted by identifying tetrahedral faces that occur
only once. Patches grow through triangle-edge adjacency from a randomly selected
surface vertex and an incident triangle, until their accumulated reference area
reaches the requested target. These are connected surface patches, not geometric
circles. The area overshoot is less than the area of the final added triangle.

For independent draws $U\sim\mathcal{U}[0,1)$, the current sampler uses
\begin{align}
 r_D &= 0.025\,\mathrm{m}+0.05\,U L_{\max}, & A_D^{\mathrm{target}}&=3\pi r_D^2,\\
 r_p &= (0.01+0.06U)\,\mathrm{m}, & A_p^{\mathrm{target}}&=\pi r_p^2.
\end{align}
Here $L_{\max}$ is the largest axis-aligned extent of the reference mesh. The
first term in $r_D$ is an absolute length; only its second term scales with liver
size. The force-patch radius parameter ranges from 1 to 7\,cm. The single-sample
default samples one to three force patches; the batch default samples two to three.

Force patches exclude the fixed patch and a conservative 2\,cm Euclidean
clearance around it. A triangle is eligible only if its centroid-to-fixed-surface
distance minus its maximum centroid-to-vertex distance is at least 2\,cm.
Force patches cannot share triangles, although they may share edges or vertices.
An insufficient eligible connected area produces a preparation failure without
resampling, shrinking a patch, or reducing a force.

All three displacement components are zero on the fixed patch. For each loaded
patch, we sample a nominal force magnitude and direction as
\begin{equation}
 q_p=s(1-U_p^2)\,\mathrm{N},\qquad
 \vect{d}_p=\frac{\vect{v}_p}{\|\vect{v}_p\|_2},\qquad
 (v_p)_i\sim\mathcal{U}[-0.5,0.5).
\end{equation}
The direction distribution is induced by normalizing a vector sampled in a cube;
it is not exactly uniform on the sphere. The force-factor parameter $s$ is an
upper bound on each nominal force draw, not the force assigned to every patch.
Using the actual grown reference area $A_{p,0}$, the prescribed Cartesian traction is
\begin{equation}
 \vect{t}_p=\frac{q_p\vect{d}_p}{A_{p,0}}.
 \label{eq:traction}
\end{equation}
Its components are rounded to six decimal places before writing the SIF.

\textbf{Area convention in the current Elmer case.}
The SIF specifies \texttt{Force 1}, \texttt{Force 2}, and \texttt{Force 3}, and
does not enable \texttt{Pseudo-Traction}. In the installed Elmer implementation,
these are Cartesian tractions per unit \emph{deformed} surface area. Consequently,
the prescribed components remain constant but the total patch force varies with
its deformed area. Equation~\eqref{eq:traction} normalizes the force at the reference
configuration. A constant force obtained by integrating over reference area would
be a different loading convention and must be selected explicitly in a replacement.

\section{Material and equilibrium}
We model the liver as homogeneous, isotropic, compressible neo-Hookean material.
Young's modulus is specified by the caller or sampled uniformly from 500 to
5000\,Pa. The default Poisson ratio is $\nu=0.35$.
The Lam\'e parameters are
\begin{equation}
 \mu=\frac{E}{2(1+\nu)},\qquad
 \lambda=\frac{E\nu}{(1+\nu)(1-2\nu)}.
\end{equation}
Define the deformation gradient, volume ratio, and right Cauchy--Green tensor by
\begin{equation}
 \vect{F}=\vect{I}+\nabla_X\vect{u},\qquad
 J=\det\vect{F}>0,\qquad \vect{C}=\vect{F}^{T}\vect{F}.
\end{equation}
The strain-energy density corresponding to the stress law in Elmer's
\texttt{NeoHookeanLocalMatrix} is
\begin{equation}
 W(\vect{F})=
 \frac{\mu}{2}\bigl(\operatorname{tr}\vect{C}-3\bigr)-\mu\ln J
 +\frac{\lambda}{4}\bigl(J^2-1-2\ln J\bigr).
 \label{eq:energy}
\end{equation}
Its first Piola--Kirchhoff stress is
\begin{equation}
 \vect{P}=\frac{\partial W}{\partial\vect{F}}
 =\mu(\vect{F}-\vect{F}^{-T})
   +\frac{\lambda}{2}(J^2-1)\vect{F}^{-T}.
\end{equation}
Several distinct compressible material laws are called neo-Hookean. Replacing
Equation~\eqref{eq:energy} with a different volumetric energy can change the
deformation even when $E$ and $\nu$ are unchanged.

Equilibrium requires zero force imbalance for all admissible virtual
displacements $\vect{w}$ that vanish on $\Gamma_D$:
\begin{equation}
 \int_{\Omega_0}\vect{P}:\nabla_X\vect{w}\,dV_0
 -\sum_p\int_{\Gamma_{p,0}}\vect{w}\cdot\vect{t}_p\,j_s\,dA_0=0,
 \qquad
 j_s=J\|\vect{F}^{-T}\vect{N}\|_2.
 \label{eq:weak}
\end{equation}
Here $\vect{N}$ is the reference surface normal and $j_s$ is the deformed-to-reference
area ratio. Setting $j_s=1$ in the load integral instead describes prescribed
reference-area traction. The unforced boundary has zero prescribed traction but
is free to move.

\section{Tetrahedral discretization}
Within each tetrahedron $e$, linear shape functions $N_a$ interpolate displacement:
\begin{equation}
 \vect{u}_h(\vect{X})=\sum_{a=1}^{4}N_a(\vect{X})\vect{u}_a,
 \qquad
 \vect{F}_e=\vect{I}+\sum_{a=1}^{4}
       \vect{u}_a\otimes\nabla_X N_a.
\end{equation}
The shape-function gradients and reference volume $V_{e,0}$ are constant in a
linear tetrahedron and can be precomputed. The element's internal nodal force is
\begin{equation}
 \vect{f}^{\mathrm{int}}_{e,a}
 =V_{e,0}\vect{P}_e\nabla_X N_a.
\end{equation}
For uniform traction on a boundary triangle of deformed area $A_f(\vect{u})$,
each of its three vertices receives the elemental contribution
\begin{equation}
 \vect{f}^{\mathrm{ext}}_{f,a}
 =\frac{A_f(\vect{u})}{3}\vect{t}_p.
\end{equation}
Contributions are accumulated at shared vertices. Boundary group IDs select
the prescribed traction; the native boundary file itself stores no force values.

Assembly gives the global residual
\begin{equation}
 \vect{r}(\vect{u})=
 \vect{f}^{\mathrm{int}}(\vect{u})-
 \vect{f}^{\mathrm{ext}}(\vect{u}).
\end{equation}
The equilibrium equations are solved on unconstrained degrees of freedom.
Nonzero residuals at fixed vertices represent support reactions and must not be
treated as unbalanced free-node forces.

\section{Nonlinear solution and line search}
A Newton formulation computes an update from
\begin{equation}
 \vect{K}_k\,\Delta\vect{u}_k=-\vect{r}(\vect{u}_k),\qquad
 \vect{K}_k=\frac{\partial\vect{r}}{\partial\vect{u}}
             \bigg|_{\vect{u}_k},\qquad
 \vect{u}_{k+1}=\vect{u}_k+\alpha\Delta\vect{u}_k.
 \label{eq:newton}
\end{equation}
The tangent changes because material response and geometry depend on the current
displacements. For the current-area traction convention, a fully consistent
tangent also contains the derivative of the external nodal load. Elmer's boundary
routine states that these load dependencies are evaluated using the previous
iterate; an independently implemented consistent Newton tangent therefore need
not reproduce Elmer's iteration history, even for the same equilibrium problem.

The current default solves each linearized system using CG with Cholesky2
preconditioning, a maximum of 500 linear iterations, and a linear tolerance of
$10^{-7}$. CG requires a symmetric positive-definite system; this property is
not guaranteed for every large-deformation tangent. GMRES/ILU2 and
BiCGStabl/ILU2 templates are also available. Replacing the linear solver does
not change the intended equilibrium equations.

Elmer's default line-search cost is the norm of the assembled algebraic residual.
Writing this cost as $c$, the configured search accepts a trial when
\begin{equation}
 c(\vect{u}_k+\alpha\Delta\vect{u}_k)
 \leq (1-\eta\alpha)c(\vect{u}_k).
 \label{eq:armijo}
\end{equation}
The full step starts at $\alpha=1$ and rejected steps are halved. The current
templates inherit Elmer's defaults of $\eta=0.5$ and ten trials. A separate replay
experiment used $\eta=0.1$, requiring a 10\% cost reduction for a full step instead
of 50\%. A trial reassembles and evaluates the residual at a new displacement;
it does not perform another CG solve. The same update direction is reused until
the trial is accepted or the search fails.

\section{Stopping and validation}
The current line-search path uses the change in the scalar displacement norm as
its nonlinear stopping measure. For $M$ displacement components,
\begin{align}
 n_k&=\sqrt{\frac{1}{M}\sum_{i=1}^{M}u_{k,i}^2},\\
 \delta_k&=\frac{2|n_k-n_{k-1}|}{n_k+n_{k-1}}<10^{-4}.
 \label{eq:stop}
\end{align}
The zero-denominator case is handled separately by Elmer. This dimensionless
threshold corresponds approximately to a change below 0.01\% in the overall
RMS displacement. It is neither a per-vertex displacement threshold nor a force
equilibrium tolerance. In particular, the difference between two norms is not
the norm of the difference between two displacement fields.

The script accepts a simulation only when the solver exits successfully,
reports no Elmer error, produces its result file, and has a finite final logged
nonlinear change below the SIF tolerance. It additionally requires finite
coordinates and displacements and no tetrahedron with zero or reversed signed
volume relative to its reference state. These checks do not independently
establish a specified final force-residual tolerance or exclude surface
self-intersection.

\begin{tabularx}{\linewidth}{@{}lX@{}}
\toprule
Failure category & Interpretation \\
\midrule
Preparation failure & An eligible connected surface region cannot supply the sampled patch area. \\
Line-search failure & None of the allowed trial step lengths satisfies sufficient reduction. \\
Linear divergence & The linear residual exceeds the solver's divergence limit; this is distinct from reaching the iteration cap. \\
Nonlinear nonconvergence & The final logged change fails the required tolerance, commonly after exhausting 100 iterations. \\
Invalid geometry & Nonfinite output or a collapsed/inverted tetrahedron is detected. \\
\bottomrule
\end{tabularx}

\section{Current implementation and files}
The input VTK is read into coordinate and connectivity arrays. These directly
generate the following Elmer input files:
\begin{center}
\begin{tabularx}{\linewidth}{@{}lX@{}}
\toprule
File & Content \\
\midrule
\path{mesh.nodes} & Node IDs and reference coordinates. \\
\path{mesh.elements} & Tetrahedral connectivity and body IDs. \\
\path{mesh.header} & Counts and element types. \\
\path{mesh.boundary} & Surface triangles, parent tetrahedra, and boundary group IDs. \\
\path{case.sif} & Material properties, group boundary values, and solver options. \\
\bottomrule
\end{tabularx}
\end{center}
Group 1 is free, group 2 is fixed, and groups 3 onward represent individual load
patches. Elmer reads these native files; it does not read the source VTK directly.
There are no intermediate \path{volume.vtk} or \path{simulation.vtk} files.

The default template specifies one steady-state solve, 100 nonlinear iterations,
and enabled line search. It also records Newton-after-iterations 50,
Newton-after-tolerance $10^{-3}$, and relaxation factor 1. These keywords should
not be interpreted as proof of 50 initial Picard steps: the installed
neo-Hookean element routine explicitly constructs a Newton material tangent.
No load ramp is currently configured.

\path{manifest.json} records the arguments, sampled material, patch seeds and
areas, prescribed tractions, full SIF, timing, and status. Elmer writes
\path{result/case_t0001.vtu}; its coordinates are already deformed. On acceptance,
the script writes \path{deformed.vtk}. Displacements must not be added to those
deformed coordinates a second time. Batch processing additionally records
per-sample metrics and optional rigid-registration TRE.

For example, an input already expressed in metres can be processed by:
\begin{lstlisting}
pixi run --manifest-path baselines/Volume2SurfaceCNN/pixi.toml python \
  simulate_deformation/main.py liver.vtk outputs/example_deformation \
  --seed 1 --force-factor 3 --poisson-ratio 0.35 \
  --youngs-modulus 3000 --min-force-patches 2 --template cg
\end{lstlisting}
The output directory must be new. The line-search replay changed only the
following SIF setting; it did not modify the default templates:
\begin{lstlisting}
Nonlinear System Linesearch Limit = 0.1
\end{lstlisting}

\section{Example custom-solver outline}
A replacement can retain the existing mesh reader, surface extraction, patch
sampler, manifests, and VTK output. It needs a new material/element kernel,
sparse assembly, and nonlinear driver. Existing sparse linear algebra libraries
can solve the linear systems; implementing CG or a sparse factorization from
scratch is unnecessary.

\begin{lstlisting}
read reference vertices and tetrahedra in metres
extract surface; sample fixed and traction patches
precompute tetrahedral gradients, volumes, and assembly indices
identify fixed and free displacement components
u = 0

for load_fraction in [0.1, 0.2, ..., 1.0]:
    scale prescribed tractions by load_fraction
    for nonlinear_iteration in range(max_iterations):
        compute element F, det(F), stresses, and internal forces
        assemble external forces using the chosen area convention
        assemble residual on free components
        check explicitly defined residual convergence
        assemble consistent tangent; solve K * du = -residual

        alpha = 1
        for trial in range(max_line_search_trials):
            trial_u = u + alpha * du
            reject trial if any tetrahedron has det(F) <= 0
            recompute trial residual
            if sufficient reduction is satisfied:
                accept trial_u; break
            alpha *= 0.5
        record residual, update size, alpha, and iteration counts
        stop and record failure if no trial is accepted
    stop if the load increment did not converge

validate geometry and final force residual at load_fraction = 1
write deformed VTK and manifest
\end{lstlisting}

This proposed continuation procedure retains the same mesh, material, patches,
and final tractions. A partial-load solution must not be reported as a successful
full-load deformation. Incremental loading is also discussed for the FEM data
generation in U-Mesh~\cite{mendizabal}.

For an independent solver, we would explicitly monitor the free-node residual,
for example $\|\vect{r}_{\mathrm{free}}\|_2/
\|\vect{f}^{\mathrm{ext}}_{\mathrm{free}}\|_2$ when the denominator is nonzero,
together with displacement-update size and element Jacobians. Its tolerance
would need validation; it must not be assigned $10^{-4}$ on the assumption that
it is equivalent to Equation~\eqref{eq:stop}.

Verification should include zero-load and rigid-motion material checks,
finite-difference checks of element forces and tangent derivatives, a
small-deformation comparison against linear elasticity, and matched comparisons
with converged Elmer cases. Matching the energy, traction convention, and
constraints is necessary before attributing differences to solver accuracy.

\section*{Implementation references}
The implementation described here is based on the repository state inspected on
18 September 2026. The main source locations are:
\begin{itemize}
  \item \path{simulate_deformation/boundaries.py}: patch and force sampling.
  \item \path{simulate_deformation/main.py}: preparation, execution, and acceptance.
  \item \path{simulate_deformation/case.sif.template}: default solver settings.
  \item \path{simulate_deformation/elmer_mesh.py}: native mesh writing.
  \item \path{simulate_deformation/mesh_utils.py}: geometry validation.
  \item Elmer \path{fem/src/modules/ElasticSolve.F90}:\\
    \texttt{NeoHookeanLocalMatrix} and \texttt{LocalBoundaryMatrix}.
  \item Elmer \path{fem/src/SolverUtils.F90}:
    \texttt{CheckStepSize} and \texttt{ComputeNorm}.
\end{itemize}
The inspected Elmer installation reports version 26.2, revision
\texttt{a19504a}; its source is stored under
\path{baselines/Volume2SurfaceCNN/.elmer-source/}. Its formulations are documented
in the Elmer Models Manual~\cite{elmer}. LIBR's
\path{baselines/LIBR/src/libr/fem.py} supplies reusable linear-tetrahedron geometry
and assembly patterns, but its linear material stiffness is not a neo-Hookean solver.

\begin{thebibliography}{9}
\bibitem{elmer}
CSC--IT Center for Science. \emph{Elmer Models Manual}, chapter on finite elasticity.
\url{https://www.nic.funet.fi/index/elmer/doc/ElmerModelsManual.pdf}.
The installed source is the authority for the implementation details above.

\bibitem{mendizabal}
A. Mendizabal, P. M\'arquez-Neila, and S. Cotin.
Simulation of hyperelastic materials in real-time using deep learning.
\emph{Medical Image Analysis}, 59:101569, 2020.
\url{https://arxiv.org/abs/1904.06197}.
\end{thebibliography}
\end{document}
